\section{Experimental Setup}

This section details the experimental configuration for validating Malak Platform on the CIFAR-10 image classification benchmark.

\subsection{Application: CIFAR-10 Image Classification}

\subsubsection{Task Description}
Object classification in natural images for edge device deployment. This benchmark validates the platform's ability to train, compress, and deploy vision models efficiently on resource-constrained devices.

\subsubsection{Dataset}
\begin{itemize}
\item \textbf{Source}: CIFAR-10 dataset \cite{cifar10}
\item \textbf{Size}: 50,000 training images, 10,000 test images
\item \textbf{Classes}: 10 object categories (airplane, automobile, bird, cat, deer, dog, frog, horse, ship, truck)
\item \textbf{Resolution}: 32$\times$32 RGB images
\item \textbf{Preprocessing}: Mean-variance normalization using dataset statistics
\item \textbf{Augmentation}: Random horizontal flip, random crop with 4-pixel padding
\end{itemize}

\subsubsection{Model Architecture}
\begin{itemize}
\item \textbf{Baseline}: MobileNetV2 architecture adapted for CIFAR-10
\item \textbf{Parameters}: 2,236,682 (approximately 2.2M)
\item \textbf{Modifications}: Adjusted first convolutional layer for 32$\times$32 input, modified classifier for 10 classes
\item \textbf{Model size (FP32)}: 8.76 MB
\end{itemize}

\subsubsection{Training Configuration}
\begin{verbatim}
Optimizer: SGD (lr=0.01, momentum=0.9,
           weight_decay=5e-4)
Batch size: 128
Epochs: 100
Learning rate schedule: Cosine annealing
Augmentation: Random horizontal flip,
              random crop (padding=4)
Device: CPU (Intel/AMD x86_64)
\end{verbatim}

\subsubsection{Compression Pipeline}

We evaluate the following compression strategies:

\textbf{Post-Training Quantization (PTQ)}:
\begin{itemize}
\item Dynamic INT8 quantization of linear and convolutional layers
\item No retraining required
\item Applied using PyTorch's native quantization API
\end{itemize}

\textbf{Quantization-Aware Training (QAT)}:
\begin{itemize}
\item Simulated quantization during training
\item Model adapts to reduced precision
\item Additional 20 fine-tuning epochs with reduced learning rate
\end{itemize}

\subsection{Evaluation Metrics}

\subsubsection{Accuracy Metrics}
\begin{itemize}
\item \textbf{Top-1 Accuracy}: Percentage of correctly classified test images
\item \textbf{Accuracy Degradation}: Percentage point difference from FP32 baseline
\end{itemize}

\subsubsection{Efficiency Metrics}
\begin{itemize}
\item \textbf{Model Size}: Storage required for model weights (MB)
\item \textbf{Compression Ratio}: FP32 size divided by compressed size
\item \textbf{Inference Latency}: Average milliseconds per image (averaged over 10,000 test images)
\item \textbf{Parameters}: Total trainable parameters in the model
\end{itemize}

\subsection{Hardware Platform}

Experiments were conducted on:
\begin{itemize}
\item \textbf{CPU}: x86\_64 architecture
\item \textbf{Framework}: PyTorch 2.10.0
\item \textbf{Python}: 3.11
\item \textbf{Operating System}: Linux
\end{itemize}

\subsection{Baseline Comparisons}

We evaluate Malak Platform's compression pipeline against:
\begin{itemize}
\item \textbf{FP32 Baseline}: Uncompressed full-precision model
\item \textbf{INT8 PTQ}: Post-training dynamic quantization
\item \textbf{INT8 QAT}: Quantization-aware training (estimated)
\end{itemize}

\subsection{Reproducibility}

To ensure reproducibility:
\begin{itemize}
\item \textbf{Public dataset}: CIFAR-10 is freely available and widely used
\item \textbf{Standard architecture}: MobileNetV2 is a well-established model
\item \textbf{Fixed random seed}: Set to 42 for consistent results
\item \textbf{Documented pipeline}: Complete training script provided in repository
\item \textbf{Hyperparameters}: All settings explicitly specified
\end{itemize}

Researchers can replicate experiments using:
\begin{verbatim}
git clone https://github.com/alnemari-m/malak_platform
cd malak_platform
python simple_experiment.py
\end{verbatim}
